% Source: https://artofproblemsolving.com/community/c3h1449538_diophantine_equation
% Stable identity: aops-topic-1449538-q01
We show that the solutions are exactly the pairs given by
\[
x_n+y_n\sqrt2=(3+2\sqrt2)^n,\qquad n=1,2,3,\dots,
\]
that is, $(x_1,y_1)=(3,2)$ and $x_{n+1}=3x_n+4y_n$, $y_{n+1}=2x_n+3y_n$: $(3,2),(17,12),(99,70),(577,408),\dots$

\textbf{These are solutions.} If $x+y\sqrt2=(3+2\sqrt2)^n$ with integers $x,y$, then $x-y\sqrt2=(3-2\sqrt2)^n$ (the conjugate of a product is the product of the conjugates), so
\[
x^2-2y^2=(x+y\sqrt2)(x-y\sqrt2)=\left((3+2\sqrt2)(3-2\sqrt2)\right)^n=1^n=1 .
\]
Equivalently, the recurrence preserves the equation: $(3x+4y)^2-2(2x+3y)^2=x^2-2y^2$. Both coordinates are positive.

\textbf{There are no others (descent).} Let $(x,y)$ be a solution with $y\ge1$ and put
\[
x'=3x-4y,\qquad y'=3y-2x,\qquad\text{so that}\qquad x'+y'\sqrt2=(x+y\sqrt2)(3-2\sqrt2).
\]
\begin{itemize}
\item $x'^2-2y'^2=x^2-2y^2=1$ (the same computation as above).
\item $x'>0$: $9x^2=9+18y^2>16y^2$, so $3x>4y$.
\item $y'\ge0$: $y=1$ would give $x^2=3$, impossible, so $y\ge2$; then $9y^2\ge4+8y^2=4x^2$, so $3y\ge2x$.
\item $y'<y$: this means $x>y$, true since $x^2=1+2y^2>y^2$.
\end{itemize}
So from any solution with $y\ge1$ we get a solution $(x',y')$ with $0\le y'<y$, and $(x,y)$ is recovered by $x+y\sqrt2=(x'+y'\sqrt2)(3+2\sqrt2)$. Repeating, $y$ strictly decreases until it reaches $0$, where $x^2=1$ gives the pair $(1,0)$, that is $1+0\sqrt2$. Undoing the $n$ steps, $x+y\sqrt2=(3+2\sqrt2)^n$ with $n\ge1$. $\blacksquare$
